% ============================================================
% Отчёт по исследованию ЯМ — баскетбол
% Адиатуллин Т. Р., СПбПУ, 2026
% ============================================================

\documentclass[a4paper, final]{article}

\usepackage[russian]{babel}
\usepackage{fontspec}
\setmainfont{Times New Roman}
\newfontfamily\cyrillicfonttt{Courier New}
\usepackage[14pt]{extsizes}
\usepackage[left=20mm, top=20mm, right=20mm, bottom=20mm, footskip=10mm]{geometry}
\usepackage{ragged2e}
\usepackage{indentfirst}
\usepackage{graphicx}
\usepackage{tikz}
\usepackage{caption}
\usepackage{xcolor}
\usepackage{subcaption}
\usepackage{booktabs}
\usepackage{hyperref}
\usepackage{array}
\usepackage{tabularx}
\usepackage{longtable}
\usepackage{multirow}
\usepackage{pdflscape}
\usepackage{soul}
\usepackage{enumitem}

\definecolor{link_color}{HTML}{0000EE}

\hypersetup{
    colorlinks=true,
    linkcolor=link_color,
    urlcolor=link_color,
    citecolor=link_color
}

% ============================================================
% Вспомогательные команды
% ============================================================

\newcommand{\scoretablefilled}[8]{%
    \vspace{4pt}
    \begin{table}[h!]
    \centering
    \small
    \begin{tabular}{|l|c|c|}
    \hline
    \textbf{Критерий (1--5)} & \textbf{Grok} & \textbf{Qwen} \\
    \hline
    Достоверность содержания & #1 & #2 \\
    \hline
    Полнота и структурированность & #3 & #4 \\
    \hline
    Уместность стиля & #5 & #6 \\
    \hline
    Искажения / ошибки (да/нет) & #7 & #8 \\
    \hline
    \end{tabular}
    \caption{Оценка ответов}
    \end{table}
    \vspace{2pt}
}

% ============================================================
\begin{document}

\tableofcontents
\newpage

% ============================================================
\section{Приложение: Адиатуллин Т.\,Р. — Предметная область «Баскетбол»}

Исследование проводилось на двух базовых вопросах: \textbf{пик-н-ролл} (тактика) и \textbf{тайм-аут} (правила). Каждый вопрос был сформулирован в девяти вариантах: пять функциональных стилей и четыре типа речевых помех. Все запросы задавались моделям \textbf{Grok} (xAI) и \textbf{Qwen} (Alibaba).

\newpage

% ============================================================
\subsection{Вопрос 1: Тактика — пик-н-ролл}

% ------------------------------------------------------------
\subsubsection{Эксперимент 1 — Научный стиль}

\noindent\textbf{Гипотеза.} Ожидается структурированный ответ в академическом стиле, разделённый на логические блоки: механизм выполнения, роли участников (скринер и плеймейкер) и условия эффективного применения. Лексика должна быть нейтральной и точной, без разговорных элементов. Должны быть использованы термины: заслон, скринер, ролл, свитч, дабл-тим.

\noindent\textbf{Запрос:} Дайте системное описание тактического взаимодействия «заслон-выход» (пик-н-ролл) в баскетболе: механизм выполнения, роли участников и условия эффективного применения.

\noindent\textbf{Ответ Grok:} В баскетболе тактическое взаимодействие заслон-выход или пик-н-ролл строится на взаимодействии двух игроков: один ставит заслон мячу, второй с мячом использует его для прохода или броска. Механизм начинается с того что игрок с мячом ведет дриблинг навстречу партнеру который занимает позицию сбоку от защитника и плотно ставит неподвижный заслон корпусом не двигаясь с места. Игрок с мячом проходит вплотную к заслонившему партнеру вынуждая защитника выбирать между преследованием и сменой опеки. После прохода заслоняющий разворачивается лицом к кольцу открываясь для передачи и потенциального броска или прохода под щит. Роли четко распределены: первый номер ведет атаку и решает бросок проход или передачу второй создает пространство и получает опцию на свободный бросок. Эффективность достигается при хорошей координации быстром чтении защиты точном тайминге движения и умении обоих игроков завершать атаку в случае свитча или дабл-тим.

\noindent\textbf{Ответ Qwen:} Тактическое взаимодействие заслон-выход представляет собой фундаментальный элемент баскетбольной игры, основанный на взаимодействии двух игроков для создания пространства и преимущества в атаке. Механизм выполнения начинается с того, что игрок без мяча ставит плотный неподвижный заслон партнеру с мячом, блокируя линию движения защитника атакующего игрока. Сразу после постановки заслона устанавливающий его игрок совершает резкий выход к корзине или откатывается на периферию в зависимости от тактического замысла и реакции защиты. Роли участников чётко разделены: первый игрок является создателем помех и потенциальным завершителем атаки у кольца, тогда как второй игрок с мячом должен использовать этот момент для обострения ситуации, либо прорываясь к кольцу, либо открываясь для дальнего броска. Условием эффективного применения является своевременность и качество самого заслона, который должен быть поставлен до начала движения партнера, чтобы задержать защитника.

\scoretablefilled{5}{5}{4}{5}{3}{4}{нет}{нет}

\noindent\textbf{Анализ и вывод.} Оба ответа фактически верны. Grok упоминает свитч и дабл-тим — профессиональные термины, уместные в научном контексте. Однако ответ Grok написан полностью без знаков препинания, что недопустимо для научного стиля: текст воспринимается как неструктурированный поток. Выражение «первый номер» — разговорный жаргон, несовместимый с академическим регистром. Qwen выдерживает научный стиль последовательнее: чётко разграничивает механизм, роли и условия эффективности в соответствии со структурой запроса. Единственный изъян — грамматически неловкая конструкция «критически важно также чтение защитой». Гипотеза подтверждена частично: структура ответа присутствует, однако Grok не выдержал стилистические требования.

\newpage

% ------------------------------------------------------------
\subsubsection{Эксперимент 2 — Официально-деловой стиль}

\noindent\textbf{Гипотеза.} Ожидается нейтральный, структурированный ответ в духе делового документа: без образных выражений, с чётким перечислением порядка исполнения и функций каждого игрока. Лексика должна быть сухой и точной; эмоционально окрашенные слова недопустимы.

\noindent\textbf{Запрос:} Прошу предоставить структурированное описание тактического приёма «пик-н-ролл», включая порядок его исполнения и функции каждого из задействованных игроков.

\noindent\textbf{Ответ Grok:} Тактический приём пик-н-ролл в баскетболе представляет собой взаимодействие двух игроков для создания преимущества в атаке. Один игрок устанавливает заслон для партнёра с мячом, а тот использует этот заслон для прохода или броска. Исполнение начинается с того, что игрок без мяча ставит плотный экран рядом с держателем мяча, обычно на периметре. Затем держатель мяча двигается вдоль заслона, заставляя защитника выбирать между следованием за ним или переключением. После заслона игрок, ставивший экран, открывается для передачи или откатывается под кольцо. Задерживающий мяч читает защиту и решает, атаковать корзину, отдать пас или бросить. Игрок с мячом создаёт пространство и угрозу прохода, а его партнёр обеспечивает заслон и последующую угрозу.

\noindent\textbf{Ответ Qwen:} Пик-н-ролл является фундаментальным тактическим приёмом в баскетболе, представляющим собой взаимодействие двух игроков для создания преимущества при атаке. Один игрок, называемый скринером, выставляет плотный заслон (пик) на пути защитника атакующего игрока, который владеет мячом и движется к корзине или периметру. Цель скринера — физически заблокировать движение соперника, заставив его отстать или запутаться, после чего он должен развернуться и пойти под кольцо (ролл), чтобы получить пас вблизи щита или открыть линию передачи. Атакующий с мячом использует этот момент замешательства, чтобы либо прорваться сквозь защиту к кольцу, либо совершить бросок с дистанции, если защитник пропустил его вперёд. Вторая линия защиты часто вынуждена реагировать на эти действия, создавая дополнительные возможности для передачи третьему игроку или повторного использования схемы.

\scoretablefilled{5}{5}{4}{5}{3}{3}{нет}{нет}

\noindent\textbf{Анализ и вывод.} Оба ответа фактически точны, однако ни один не выдержал официально-деловой стиль в полной мере. У Grok выражение «задерживающий мяч» звучит нестандартно, а конструкция «создаёт пространство и угрозу» ближе к спортивной журналистике. У Qwen слова «момент замешательства» и «запутаться» — эмоционально окрашенная лексика, не характерная для делового документа. Qwen при этом полнее раскрывает механизм, включая реакцию второй линии защиты. Гипотеза не подтверждена: обе модели систематически отклонялись от делового регистра.

\newpage

% ------------------------------------------------------------
\subsubsection{Эксперимент 3 — Публицистический стиль}

\noindent\textbf{Гипотеза.} Ожидается живой, динамичный ответ с риторическими элементами, образными сравнениями и объяснением актуальности приёма в современном баскетболе. Ответ должен быть убедительным и читабельным, содержать конкретные примеры или отсылки к реальной игровой практике.

\noindent\textbf{Запрос:} Пик-н-ролл — простейший приём или гениальная ловушка для защиты? Объясни, как работает эта связка и почему она до сих пор остаётся оружием номер один в НБА.

\noindent\textbf{Ответ Grok:} Пик-н-ролл это одновременно простейший и самый эффективный приём в баскетболе потому что он создаёт постоянный выбор для защиты который почти невозможно решить идеально. Один игрок ставит заслон своему партнёру с мячом тот использует его чтобы обойти защитника а заслоняющий затем разворачивается к кольцу. Защита вынуждена либо переключаться либо оставаться что открывает либо проход под кольцо либо бросок из-за дуги либо открытого третьего игрока. В НБА он остаётся оружием номер один потому что пространство площадки огромно трёхочковые броски смертельно точны а скорость и атлетизм позволяют разыгрывающим и большим идеально читать ситуацию создавая преимущество на каждом владении даже против самых умных схем обороны.

\noindent\textbf{Ответ Qwen:} Пик-н-ролл часто ошибочно считают простейшим приёмом, однако в реальности это сложная тактическая ловушка, которая постоянно эволюционирует и остаётся основным оружием НБА благодаря своей универсальности. Суть связки заключается во взаимодействии двух игроков: одного, который ставит заслон (пик) для игрока с мячом, чтобы освободить его от опеки, и другого, который после передачи бежит к кольцу (ролл). Защита оказывается перед дилеммой: либо один защитник идёт на перекрывание пути атакующему с мячом, позволяя пасовать на роллера, либо защита задерживает роллера, открывая возможность для броска или прохода с мячом. Гениальность этого элемента в том, что он вынуждает защиту совершать ошибки независимо от выбранной стратегии, будь то агрессивное давление, подстраховка или переключение между игроками. Постоянное совершенствование вариантов исполнения, включая пик-н-поп и флэш-роллы, делает этот приём незаменимым инструментом для создания лёгких очков, гарантируя его доминирование на высшем уровне баскетбола уже более века.

\scoretablefilled{5}{5}{4}{5}{4}{5}{нет}{нет}

\noindent\textbf{Анализ и вывод.} Публицистический стиль оказался наиболее органичным для обеих моделей. Grok верно схватил риторическую интонацию и объяснил современные причины эффективности приёма через пространство площадки и трёхочковую угрозу. Qwen выигрывает по полноте: упоминание пик-н-поп и флэш-роллов свидетельствует о понимании эволюции тактики. Фактическая неточность у Qwen: «другого, который после передачи бежит к кольцу» — ролл выполняется сразу после заслона, а не после передачи. Гипотеза подтверждена: оба ответа живые, динамичные, содержательные.

\newpage

% ------------------------------------------------------------
\subsubsection{Эксперимент 4 — Разговорный стиль}

\noindent\textbf{Гипотеза.} Ожидается непринуждённый ответ в бытовом тоне: короткие предложения, простые слова, отсутствие терминологической нагрузки. Модель должна «переключиться» в неформальный регистр вслед за запросом и объяснить приём так, как это делают в живом разговоре.

\noindent\textbf{Запрос:} слушай объясни мне нормально что такое пик-н-ролл в баскетболе и как его делают

\noindent\textbf{Ответ Grok:} Пик-н-ролл в баскетболе это один из самых распространенных и эффективных тактических приемов в нападении. Один игрок ставит заслон своему партнеру с мячом, а потом сразу разворачивается и идет под кольцо за передачей. Делают это так: атакующий с мячом ведет его, его партнер подходит сбоку и плотно встает перед защитником, закрывая ему путь. После контакта игрок с мячом огибает заслон и либо бросает, либо проходит к кольцу, а заслонщик в это время поворачивается лицом к кольцу и получает пас на открытый бросок или данк. Все зависит от реакции защиты.

\noindent\textbf{Ответ Qwen:} Пик-н-ролл это фундаментальная комбинация в баскетболе, состоящая из двух игроков, где один ставит заслон другому для создания преимущества. Суть действия заключается в том, что игрок без мяча подходит к защитнику того, кто владеет мячом, и встает неподвижно, блокируя его путь, чтобы тот не мог догнать партнера с мячом. После установки заслона человек, который его поставил, должен быстро развернуться и побежать к корзине, получив имя ролл, а его партнер с мячом должен использовать освободившееся пространство для прохода под кольцо или для броска. Эта тактика требует идеальной координации, понимания момента и быстрой реакции обоих участников, так как успех зависит от того, насколько точно выставлен заслон и насколько быстро разыгрывающий использует открывшуюся зону.

\scoretablefilled{5}{5}{4}{4}{5}{2}{нет}{нет}

\noindent\textbf{Анализ и вывод.} Наибольшее расхождение в уместности стиля во всём исследовании: Grok получил 5, Qwen — 2. Grok точно уловил регистр: живой темп, слово «данк», финальная фраза «всё зависит от реакции защиты» — именно так объясняют баскетбол в разговоре. Qwen полностью проигнорировал стиль запроса: «фундаментальная комбинация», «получив имя ролл» — это канцелярская лексика. Гипотеза подтверждена для Grok и не подтверждена для Qwen.

\newpage

% ------------------------------------------------------------
\subsubsection{Эксперимент 5 — Художественный стиль}

\noindent\textbf{Гипотеза.} Ожидается образный, метафоричный ответ, передающий динамику и драматизм момента. Механика приёма должна быть описана через художественные образы, но без потери фактической точности: заслон, проход, растерянность защиты — всё это должно быть передано образно, не искажая реального смысла.

\noindent\textbf{Запрос:} Два игрока, один замысел. Опиши пик-н-ролл как маленькую драму на площадке — кто ставит заслон, кто уходит в отрыв, и как защита рассыпается в этот момент.

\noindent\textbf{Ответ Grok:} На площадке два игрока с одним замыслом разыгрывают маленькую драму. Большой ставит заслон, плотно вставая перед защитником разыгрывающего, создавая стену из тела и решимости. Разыгрывающий резко уходит в отрыв, используя заслон как щит, и мчится к кольцу с мячом в руках. Защитник большого на мгновение теряется, не зная, кого опекать, а его напарник, пытаясь проскочить, натыкается на стену и отстаёт. В этот миг защита рассыпается в хаосе нерешительности, открывая путь для броска или паса, и замысел срабатывает, оставляя соперников в растерянности.

\noindent\textbf{Ответ Qwen:} В этой маленькой драме на паркете два игрока становятся главными героями, чей замысел разворачивается в едином порыве. Один из них, обладающий силой и решимостью, становится тем самым стражем, который выставляет мощный заслон, превращаясь в живую стену посреди потока игры. Его партнёр, обладая скоростью и хитростью, делает резкий рывок, имитируя уход за мячом, чтобы обмануть бдительность соперников, а затем мгновенно меняет направление к кольцу. В момент столкновения с заслоном защитник, прикрывающий ведущего, теряет равновесие или отстаёт, а тот, кто должен был помочь, оказывается слишком далеко, чтобы успеть. Пространство для броска раскрывается как сцена для финального аккорда, где защита рассыпается в хаосе недопонимания и опоздания, оставляя нападающего одного против пустой клетки.

\scoretablefilled{5}{4}{4}{5}{5}{4}{нет}{да}

\noindent\textbf{Анализ и вывод.} Grok справился лучше: компактный, образный, точный. Qwen написал развёрнуто и атмосферно, однако допустил смысловое искажение: «имитируя уход за мячом» — это описание финта, а не стандартного пик-н-ролл. Гипотеза подтверждена для Grok; для Qwen — частично, поскольку художественность достигнута за счёт фактической точности.

\newpage

% ------------------------------------------------------------
\subsubsection{Эксперимент 6 — Орфографические ошибки}

\noindent\textbf{Гипотеза.} Ожидается, что орфографические ошибки в запросе не повлияют на понимание вопроса: модели должны корректно восстановить смысл и ответить по существу. Содержание ответа должно совпадать с тем, что было бы получено при грамотном запросе.

\noindent\textbf{Запрос:} Как делается пик-н-рол в баскетболе? Объесни что делает каждый игрок и зачем эта тактека нужна?

\noindent\textbf{Ответ Grok:} В баскетболе пик-н-рол это базовая командная тактика, когда один игрок ставит заслон для партнера с мячом, а затем сам открывается под передачу. Заслоняющий игрок, обычно большой центровой или форвард, подходит к защитнику, который опекает игрока с мячом, и блокирует ему путь, создавая пространство для прохода или броска. Игрок с мячом, как правило разыгрывающий, использует этот заслон, чтобы обойти своего опекуна и либо прорваться к кольцу, либо отдать пас. После заслона первый игрок разворачивается и выходит на свободную позицию, часто на трехочковой линии или под кольцом, чтобы получить обратную передачу. Эта тактика нужна, чтобы создавать численное преимущество в атаке, нарушать оборону соперника и открывать пути для бросков или легких очков.

\noindent\textbf{Ответ Qwen:} Пик-н-ролл — это базовая атакующая тактика, в которой два игрока взаимодействуют для создания преимущества. Один из игроков, называемый скринером или поставщиком заслона, становится неподвижной преградой перед защитником игрока с мячом, который называется плеймейкером. В этот момент плеймейкер использует этот заслон, чтобы быстро проскочить мимо защиты и получить пространство для броска или прохода к кольцу. Защитник плеймейкера вынужден либо перепрыгнуть через заслон, что отнимает время, либо пропустить его, открывая путь к корзине. Скринер после постановки заслона сразу разворачивается к кольцу, становясь готовым к получению передачи и завершению атаки, если защитники переключились на него.

\scoretablefilled{4}{4}{4}{5}{3}{3}{нет}{да}

\noindent\textbf{Выявленные искажения смысла.} Grok: помеха не вызвала искажений. Содержательная неточность не связана с орфографией: «часто на трёхочковой линии» — это описание варианта пик-н-поп, а не классического пик-н-ролл. Qwen: помеха также не повлияла на понимание. Однако присутствует тактическая ошибка: «защитник вынужден либо перепрыгнуть через заслон» — в баскетболе защитник обходит заслон сверху (hedge) или снизу (go under), но не перепрыгивает его. Гипотеза подтверждена: орфографические ошибки не нарушили понимание вопроса.

\newpage

% ------------------------------------------------------------
\subsubsection{Эксперимент 7 — Дислексия}

\noindent\textbf{Гипотеза.} Ожидается, что выраженные перестановки букв (дислексия) не помешают модели восстановить смысл вопроса. Ответ должен совпадать по содержанию с ответом на грамотно сформулированный запрос, что свидетельствовало бы о высокой смысловой устойчивости.

\noindent\textbf{Запрос:} Как рабоатет пик-н-рлол в баксетболе? Что длает каждый игорк и зачем эот прёим испольузется?

\noindent\textbf{Ответ Grok:} Пик-н-рол в баскетболе это базовый приём атаки когда один игрок ставит заслон своему партнёру с мячом а потом сам получает возможность для броска или прохода. Игрок с мячом ведёт мяч и использует заслон чтобы избавиться от опеки защитника. Игрок ставящий заслон после контакта разворачивается лицом к кольцу и открывается под передачу или откатывается назад для броска. Защитники вынуждены либо менять опеку либо помогать что создаёт свободного игрока в другом месте площадки. Приём используется чтобы разрушать плотную оборону создавать численное преимущество и генерировать лёгкие очки особенно в позиционных атаках.

\noindent\textbf{Ответ Qwen:} Пик-н-ролл представляет собой классическое тактическое взаимодействие двух игроков, при котором один из них, называемый пикером, выставляет плотный блок (пик) защитнику игрока с мячом, чтобы тот не мог его сопровождать. В этот момент игрок с мячом, использующий преимущество от блока, атакует кольцо, пытаясь пройти мимо заблокированного соперника. Защитник атакующего игрока оказывается в сложной ситуации: ему нужно либо обойти блок, что занимает время и даёт преимущество атакующему, либо перепрыгнуть через него, рискуя получить фол. Параллельно сам пикер после установки блока сразу же разворачивается к корзине, чтобы получить пас от партнёра с мячом для завершения атаки или создания преимущества под кольцом.

\scoretablefilled{5}{4}{4}{4}{3}{3}{нет}{да}

\noindent\textbf{Выявленные искажения смысла.} Grok: смысловых искажений нет. Даже очень искажённые слова не помешали корректному пониманию. Qwen: вопрос понят верно, однако повторяется ошибка из эксперимента 6: «перепрыгнуть через него, рискуя получить фол» — некорректный тактический термин. Гипотеза подтверждена: дислексия не вызвала смысловых искажений.

\newpage

% ------------------------------------------------------------
\subsubsection{Эксперимент 8 — Инверсия}

\noindent\textbf{Гипотеза.} Ожидается, что нетипичный порядок слов не помешает моделям понять смысл вопроса. Ответ должен быть содержательно полным, несмотря на синтаксическую инверсию в запросе.

\noindent\textbf{Запрос:} Пик-н-ролл — как работает он? Игроки что делают каждый, и зачем тактика эта в баскетболе используется?

\noindent\textbf{Ответ Grok:} В баскетболе пик-н-ролл это базовая тактика нападения где один игрок ставит заслон для другого с мячом. Обычно большой игрок подходит и экранирует защитника разыгрывающего чтобы тот мог свободно пройти к кольцу или бросить. После заслонки большой разворачивается к кольцу открываясь для передачи и возможного броска или данка. Защитники вынуждены выбирать кого прикрывать что создает преимущество в атаке часто приводя к открытому броску или проходу под кольцо. Эта тактика используется потому что она простая эффективная и позволяет использовать скорость разыгрывающего с мощью большого игрока нарушая оборону соперника и создавая легкие очки.

\noindent\textbf{Ответ Qwen:} Пик-н-ролл — это фундаментальная тактика в баскетболе, основанная на взаимодействии двух игроков для создания численного преимущества. Суть действия заключается в том, что один игрок, обычно с мячом, выходит на дриблинг к партнёру, который выполняет роль заслона, останавливаясь перпендикулярно направлению движения атакующего. В этот момент тот, кто поставил заслон, может либо быстро развернуться и бежать к корзине, получив пас для лёгкого броска, либо остаться у линии штрафного, если защита отступит. Пик-н-ролл заставляет защитников постоянно принимать мгновенные решения, что часто приводит к ошибкам и открытым моментам для атакующей команды.

\scoretablefilled{5}{4}{4}{5}{3}{3}{нет}{да}

\noindent\textbf{Выявленные искажения смысла.} Grok: искажений нет, инверсия не повлияла на понимание. Qwen: инверсия не вызвала смысловых искажений, однако присутствуют фактические ошибки: «один игрок, обычно с мячом, выходит к партнёру» — в классическом пик-н-ролл именно игрок без мяча подходит к партнёру. «Остаться у линии штрафного» — описание пик-н-поп, а не пик-н-ролл. Гипотеза подтверждена: инверсия не нарушила понимание вопроса.

\newpage

% ------------------------------------------------------------
\subsubsection{Эксперимент 9 — Отрицание}

\noindent\textbf{Гипотеза.} Ожидается, что модели корректно разберут логику двойного отрицания и ответят о \emph{правильном} выполнении пик-н-ролл и \emph{главных} (не второстепенных) ролях участников. Запутанная синтаксическая конструкция не должна исказить смысл ответа.

\noindent\textbf{Запрос:} Не мог бы ты не отказать мне в объяснении того, как не неправильно выполняется пик-н-ролл, и какую роль не второстепенную играет каждый из участников?

\noindent\textbf{Ответ Grok:} В баскетболе пик-н-ролл выполняется правильно когда атакующий с мячом получает заслон от партнера который ставит его неподвижно и плотно к защитнику а затем сразу открывается для продолжения атаки. Игрок с мячом не должен бежать в заслон слишком рано или слишком поздно чтобы не потерять темп и не дать защитнику проскочить. Заслоняющий не должен двигаться во время контакта иначе это будет нарушение. После заслона он сразу разворачивается лицом к кольцу и готовится либо к пасу на открытый бросок либо к добиванию если мяч вернется. Остальные партнеры растягивают площадку чтобы не мешать и создавать угрозу на периметре.

\noindent\textbf{Ответ Qwen:} Пик-н-ролл часто выполняется неправильно из-за отсутствия синхронности между двумя основными участниками, а именно скринером и плеймейкером. Скринер совершает ошибку, если не создаёт плотный контакт с защитником или уходит на периметр слишком рано, лишая партнёра возможности получить свободный бросок. Плеймейкер может не прочитать ситуацию и не использовать экран эффективно, просто отводя мяч в сторону вместо того чтобы пройти сквозь него. Если один из игроков действует независимо от другого, вся конструкция рушится, так как успех этой комбинации полностью зависит от взаимопонимания и способности читать реакцию защиты в реальном времени.

\scoretablefilled{5}{5}{5}{5}{3}{3}{нет}{нет}

\noindent\textbf{Выявленные искажения смысла.} Grok: двойное отрицание понято верно. Дополнительное достоинство: единственный ответ в исследовании, где упомянута роль остальных трёх игроков (растягивают площадку). Qwen: двойное отрицание также понято верно. Модель описала типичные ошибки исполнения — оригинальный и корректный подход к вопросу. Гипотеза подтверждена для обеих моделей.

\newpage

% ============================================================
\subsection{Вопрос 2: Правила — тайм-аут}

% ------------------------------------------------------------
\subsubsection{Эксперимент 10 — Научный стиль}

\noindent\textbf{Гипотеза.} Ожидается точное и структурированное изложение регламента ФИБА: 2 тайм-аута в первой половине, 3 во второй, 1 в каждом овертайме, продолжительность 1 минута, право запроса — только у главного тренера. Ответ должен быть оформлен в академическом тоне с чёткими числовыми данными.

\noindent\textbf{Запрос:} Опишите регламентацию тайм-аутов в баскетболе согласно правилам ФИБА: количество, продолжительность, порядок запроса и ограничения по времени матча.

\noindent\textbf{Ответ Grok:} В баскетболе по правилам ФИБА каждая команда имеет право на два тайм-аута по одной минуте в каждой половине матча и один дополнительный тайм-аут по одной минуте в каждом овертайме. Запрос на тайм-аут может сделать только тренер или его помощник через секретаря или официальное лицо матча когда мяч мёртвый и часы остановлены либо когда команда владеет мячом. Тайм-аут не может быть запрошен в последние две минуты четверти или овертайма если команда только что набрала очки. Общая продолжительность тайм-аута составляет одну минуту после чего игра возобновляется.

\noindent\textbf{Ответ Qwen:} Согласно правилам ФИБА в каждом матче каждая команда имеет право запросить два тайм-аута продолжительностью одна минута во второй половине игры и один тайм-аут той же длительности в овертайме, при этом количество запрашиваемых тайм-аутов в первой половине ограничено двумя минутами суммарного времени. Тайм-аут может быть запрошен только тренером или помощником тренера до того, как мяч станет живым после свистка судьи. В случае если команда не использовала все разрешённые тайм-ауты в первом периоде, они переносятся на вторую половину матча, но не могут быть перенесены в овертайм.

\scoretablefilled{3}{2}{3}{3}{4}{4}{нет}{да}

\noindent\textbf{Анализ и вывод.} Оба ответа содержат фактические ошибки. По правилам ФИБА — 2+3+1, только главный тренер. Grok указал «два в каждой половине» — ошибка для второй половины. Qwen запутался сильнее: «ограничено двумя минутами суммарного времени» — несуществующая норма; утверждение о переносе тайм-аутов — неверно. Гипотеза не подтверждена: обе модели не воспроизвели точные регламентные данные ФИБА.

\newpage

% ------------------------------------------------------------
\subsubsection{Эксперимент 11 — Официально-деловой стиль}

\noindent\textbf{Гипотеза.} Ожидается чёткое и нейтральное изложение порядка использования тайм-аутов: кто вправе запрашивать (только главный тренер), в какой момент (при мёртвом мяче), допустимое количество (2+3+1 по ФИБА). Ответ должен быть сформулирован как деловая справка без лишних подробностей.

\noindent\textbf{Запрос:} Прошу изложить действующий порядок использования тайм-аутов в баскетболе: кто вправе их запрашивать, в какой момент игры и каково допустимое количество.

\noindent\textbf{Ответ Grok:} В баскетболе тайм-ауты запрашивают только главный тренер или его помощник, игроки на площадке не имеют права. Запрос возможен в любой момент остановки игры или когда мяч мёртвый, но не во время активной атаки. В НБА и ФИБА каждая команда получает по два тайм-аута в первой половине, по три во второй и по одному в каждом овертайме. Общее количество строго лимитировано правилами, превышение не допускается.

\noindent\textbf{Ответ Qwen:} В баскетболе право на запрос тайм-аута имеют только главные тренеры команд, которые могут делать это во время остановок игры или непосредственно перед вводом мяча в игру, при этом игроки не имеют права запрашивать паузы самостоятельно. В большинстве официальных турниров, включая правила ФИБА, команда получает пять тайм-аутов за всю игру, которые можно использовать в любой четверти. Остановка игры для тайм-аута длится ровно одну минуту, в течение которой тренер может проводить инструктаж игроков и менять тактику.

\scoretablefilled{4}{2}{4}{3}{3}{3}{нет}{да}

\noindent\textbf{Анализ и вывод.} Grok точнее: распределение 2+3+1 верное, однако упоминание «помощника тренера» противоречит ФИБА. Qwen верно назвал общее число (5), но не дал разбивки по половинам и добавил несуществующее ограничение «нельзя брать более двух подряд». Гипотеза подтверждена частично: Grok ближе к ожидаемому, Qwen значительно отклонился от реального регламента.

\newpage

% ------------------------------------------------------------
\subsubsection{Эксперимент 12 — Публицистический стиль}

\noindent\textbf{Гипотеза.} Ожидается живой, убедительный ответ о стратегическом значении тайм-аута в баскетболе: когда его берут, почему это тренерское искусство. Ответ должен быть написан в жанре спортивной публицистики — с динамикой и образностью, сфокусированный именно на баскетболе.

\noindent\textbf{Запрос:} Тайм-аут — это минута, которая может решить судьбу матча. Расскажи, когда и как команды берут тайм-аут и почему правильный выбор момента — это настоящее тренерское искусство.

\noindent\textbf{Ответ Grok:} Тайм-аут в спортивных матчах это короткая пауза обычно минута которую команда берет по просьбе тренера чтобы перестроить игру дать указания или успокоить эмоции. В баскетболе команды имеют ограниченное число тайм-аутов за период и берут их при отставании для мотивации при серии промахов или перед важным броском когда нужно обсудить тактику. Правильный момент это настоящее искусство потому что тайм-аут может сломать ритм соперника если тот на подъеме или наоборот дать ему передышку и потерять преимущество. Тренер должен чувствовать игру интуитивно оценивать усталость счет давление трибун и психологию чтобы пауза принесла максимум пользы.

\noindent\textbf{Ответ Qwen:} Тайм-аут действительно является мощным инструментом, позволяющим тренеру прервать ход игры и внести коррективы в тактику. Команды обычно берут паузу, когда соперник начинает набирать темп и отрываться в счёте, чтобы остановить его momentum, или когда своя команда проигрывает и нуждается в психологической поддержке. Правильный выбор момента превращается в настоящее искусство, поскольку слишком ранняя пауза может выдать все карты заранее, а запоздалая — уже не успеть повлиять на исход эпизода. Тренер должен чувствовать настроение игроков, читать игру противника и понимать, что именно сейчас нужно команде: успокоить ли их или добавить агрессии. Неверное решение может стоить победы, тогда как грамотная остановка способна переломить ход матча.

\scoretablefilled{4}{5}{3}{5}{3}{5}{нет}{нет}

\noindent\textbf{Анализ и вывод.} Qwen значительно лучше справился с публицистическим жанром: противопоставление «слишком ранняя — слишком поздняя», слово «momentum», нарастающее напряжение к финалу — признаки качественного публицистического текста. Grok уклонился от конкретики баскетбола и допустил фактическую ошибку: «по просьбе тренера или игрока» — игроки в ФИБА не имеют права запрашивать тайм-аут. Гипотеза подтверждена для Qwen, частично для Grok.

\newpage

% ------------------------------------------------------------
\subsubsection{Эксперимент 13 — Разговорный стиль}

\noindent\textbf{Гипотеза.} Ожидается простой и короткий ответ в бытовом тоне: модели должны объяснить количество и момент тайм-аутов так, как это делают в живом разговоре, без излишней формализации.

\noindent\textbf{Запрос:} эй а когда в баскетболе можно взять тайм-аут и сколько их вообще даётся на игру?

\noindent\textbf{Ответ Grok:} В баскетболе правила на тайм-ауты зависят от лиги но в стандартном ФИБА и НБА тайм-аут можно брать почти в любой момент когда мяч мёртвый или твоя команда владеет им перед броском. В игре обычно даётся по два тайм-аута на каждую половину в НБА или по одному в четверти плюс дополнительные в овертайме. Общее количество на матч варьируется но в среднем команда имеет около пяти-шести коротких тайм-аутов за всю игру. Тренер сигнализирует об этом судье.

\noindent\textbf{Ответ Qwen:} В баскетболе правила предоставления тайм-аутов зависят от лиги, но в профессиональном матче НБА каждая команда имеет право на семь коротких и два длинных тайм-аута за игру, при этом короткие длятся две минуты, а длинные семь минут. В международном баскетболе по правилам ФИБА команды получают пять тайм-аутов на игру длительностью одна минута каждый, причём первые три можно использовать только до конца третьей четверти, а оставшиеся два — после начала четвёртой. В студенческом баскетболе США NCAA также действуют свои особенности, где количество доступных остановок может варьироваться.

\scoretablefilled{3}{2}{3}{3}{5}{2}{нет}{да}

\noindent\textbf{Анализ и вывод.} Максимальное расхождение по стилю в блоке о тайм-аутах: Grok — 5, Qwen — 2. Grok выдержал разговорный регистр и дал приемлемо точный ответ. Qwen ответил формально и тяжело, а содержание грубо ошибочно: деление тайм-аутов НБА на «семь коротких и два длинных» продолжительностью «две и семь минут» — несуществующая схема. Гипотеза подтверждена для Grok; для Qwen — не подтверждена ни по стилю, ни по содержанию.

\newpage

% ------------------------------------------------------------
\subsubsection{Эксперимент 14 — Художественный стиль}

\noindent\textbf{Гипотеза.} Ожидается образное и атмосферное описание момента тайм-аута: напряжение, пауза, действия на площадке. Ответ должен создавать живую картину происходящего, сохраняя при этом фактическую корректность относительно самого понятия тайм-аута.

\noindent\textbf{Запрос:} Секунды тают, счёт равный, тренер поднимает руку — тайм-аут. Опиши этот момент: что это такое, когда его берут и что происходит на площадке в эту паузу.

\noindent\textbf{Ответ Grok:} В баскетболе секунды тают при равном счёде, тренер поднимает руку и берёт тайм-аут — это короткая запланированная пауза в игре, которую команда использует для передышки и корректировки. Его берут, чтобы остановить время, когда нужно обсудить тактику, успокоить игроков, дать указания или подготовить ключевой бросок в концовке матча. На площадке в эту минуту игроки собираются у скамейки, тренер быстро объясняет план, судьи фиксируют остановку часов, а зрители шумят в ожидании возобновления. Атмосфера накаляется, секунды паузы кажутся решающими для исхода поединка.

\noindent\textbf{Ответ Qwen:} Тайм-аут — это запланированная или экстренная пауза в игре, которую тренер запрашивает для остановки хода событий и обсуждения тактики. В этот момент секунды тают, но игра замирает: игроки направляются к скамейке запасных, где их окружает тренер. Внутри команды происходит бурное обсуждение: тренер объясняет ошибки прошлого отрезка, рисует новые схемы перемещений и мотивирует игроков психологически. Игроки слушают, задают вопросы, пьют воду и затирают пот со лба, готовясь вернуться в бой с обновлённой энергией. Эта минута тишины и сосредоточенности часто становится переломной, определяющей исход всего матча.

\scoretablefilled{4}{4}{3}{5}{4}{5}{нет}{нет}

\noindent\textbf{Анализ и вывод.} Qwen создал более насыщенную картину: детали «пьют воду», «затирают пот», «рисует схемы» — конкретные художественные образы. Grok точнее фактически (упомянул судей и зрителей), но текст нейтральнее. Небольшое несоответствие у Qwen: «направляются к скамейке» — в баскетболе игроки собираются у боковой линии, не покидая площадку. Гипотеза подтверждена для обеих моделей.

\newpage

% ------------------------------------------------------------
\subsubsection{Эксперимент 15 — Орфографические ошибки}

\noindent\textbf{Гипотеза.} Ожидается, что орфографические ошибки не помешают пониманию вопроса о количестве и моменте тайм-аутов. Ответ должен совпадать по содержанию с ответом на грамотно сформулированный аналогичный запрос.

\noindent\textbf{Запрос:} Сколько тайм-аутов можно взять в баскетболе и кагда именно их разришается брать по правилам?

\noindent\textbf{Ответ Grok:} В баскетболе по стандартным правилам ФИБА каждая команда может взять два тайм-аута в первой половине и три во второй половине матча, всего не более пяти за игру. В овертайме разрешается один дополнительный тайм-аут. Тайм-ауты можно брать только когда мяч мёртвый, то есть после свистка судьи, при остановке игры или после забитого мяча, но не во время активной атаки с мячом в игре. Тренер должен попросить тайм-аут у судьи в разрешённый момент, и он длится обычно одну минуту.

\noindent\textbf{Ответ Qwen:} В баскетболе в большинстве профессиональных игр, включая НБА и Евролигу, команда имеет право взять пять тайм-аутов за основную игру. В ФИБА команды могут запросить пять тайм-аутов в основное время игры, но при этом запрещено брать более двух подряд, а также запрещается использовать их в четвёртой четверти после того, как на табло останется менее двух минут. Тайм-ауты разрешается брать в любой момент, когда мяч находится в состоянии мёртвого мяча, то есть после остановки игры судьёй или по завершении розыгрыша.

\scoretablefilled{4}{2}{4}{3}{3}{3}{нет}{да}

\noindent\textbf{Выявленные искажения смысла.} Grok: орфографические ошибки не вызвали искажений. Это наиболее точный ответ Grok по теме тайм-аутов: 2+3+1 указано верно. Qwen: помеха не повлияла на понимание, однако ответ содержит выдуманные ограничения: «запрещено брать более двух подряд» и «запрещается в четвёртой четверти после двух минут» — таких норм в ФИБА не существует. Гипотеза подтверждена для Grok; для Qwen — содержание ошибочно вне зависимости от помехи.

\newpage

% ------------------------------------------------------------
\subsubsection{Эксперимент 16 — Дислексия}

\noindent\textbf{Гипотеза.} Ожидается, что сильно искажённый текст с перестановками букв будет распознан корректно, и ответ будет содержательно верным — как при грамотном запросе о количестве тайм-аутов и условиях их использования.

\noindent\textbf{Запрос:} Скольок тайм-аоутов дёатся в баксетболе и когад именон их мжоно взять по павирлам игры?

\noindent\textbf{Ответ Grok:} В баскетболе ФИБА каждая команда имеет право на два тайм-аута по одной минуте в каждой половине матча плюс один дополнительный тайм-аут в овертайме если он потребуется. В НБА правила немного другие каждая команда может взять до семи тайм-аутов за игру с ограничениями по периодам и возможностью брать их только в определённые моменты когда мяч мёртвый. Тайм-ауты можно брать только когда мяч не в игре после свистка судьи или после забитого мяча до вбрасывания.

\noindent\textbf{Ответ Qwen:} В баскетболе в НБА у команд есть 7 тайм-аутов за игру с лимитом не более трёх в четвёртой четверти, а в ФИБА правила отличаются: у команд по пять тайм-аутов за игру, из которых три доступны во второй половине, причём каждый тайм-аут длится одну минуту. Тренер должен подать сигнал арбитру о желании взять паузу, пока игра остановлена, и количество оставшихся тайм-аутов строго контролируется официальными документами для предотвращения злоупотреблений.

\scoretablefilled{3}{2}{3}{3}{3}{3}{нет}{да}

\noindent\textbf{Выявленные искажения смысла.} Grok: дислексия не вызвала смысловых искажений. Фактическая ошибка не связана с помехой: «два в каждой половине» — повторяется системная ошибка. Qwen: вопрос понят верно. Примечательно, что Qwen впервые указал «три во второй половине» — более точно, чем в предыдущих промптах. Гипотеза подтверждена: дислексия не нарушила понимание вопроса.

\newpage

% ------------------------------------------------------------
\subsubsection{Эксперимент 17 — Инверсия}

\noindent\textbf{Гипотеза.} Ожидается, что нетипичный порядок слов не нарушит понимание вопроса о том, когда берут тайм-аут, сколько их разрешается и кто может запрашивать. Ответ должен быть содержательно корректным.

\noindent\textbf{Запрос:} Тайм-аут в баскетболе — когда берут его? Сколько раз за игру разрешается и кто запрашивать его может?

\noindent\textbf{Ответ Grok:} Тайм-аут в баскетболе берут обычно когда нужно остановить игру для обсуждения тактики дать указания игрокам или чтобы сбить ритм соперника особенно при важных моментах матча или после неудачной серии бросков. В стандартных правилах ФИБА каждая команда может запросить до двух тайм-аутов в первой и второй четверти и до трёх в третьей и четвёртой плюс дополнительные в овертайме. Запрашивать тайм-аут могут только тренеры или игроки на площадке по сигналу судье.

\noindent\textbf{Ответ Qwen:} Тайм-аут в баскетболе берётся по инициативе тренера, когда мяч находится у команды на подаче или после остановки игры судьями. В НБА разрешается брать до семи тайм-аутов за игру, при этом каждый тренер может использовать не более четырёх в первой половине и трёх во второй. В международном баскетболе ФИБА командам даётся три тайм-аута за основную часть матча и один дополнительный в каждом овертайме, причём их можно переносить между периодами. Запрашивать паузу может только главный тренер или капитан команды, который должен подойти к секретарю с просьбой об остановке времени.

\scoretablefilled{2}{2}{3}{3}{3}{3}{да}{да}

\noindent\textbf{Выявленные искажения смысла.} Grok: инверсия спровоцировала фактическое искажение — тайм-ауты распределены поквартально («до двух в первой и второй четверти, до трёх в третьей и четвёртой»), тогда как ФИБА делит по половинам. «Игроки на площадке» не имеют права запрашивать тайм-аут. Qwen: «командам даётся три тайм-аута за основную часть матча» — грубая ошибка (по ФИБА их пять). «Капитан» не имеет права запрашивать тайм-аут. «Их можно переносить» — неверно. Это единственный эксперимент, где обе модели получили «да» по критерию искажений. Гипотеза не подтверждена: инверсия в сочетании с числовыми регламентными данными вызвала ошибки у обеих моделей.

\newpage

% ------------------------------------------------------------
\subsubsection{Эксперимент 18 — Отрицание}

\noindent\textbf{Гипотеза.} Ожидается, что модели корректно разберут логику двойного отрицания и ответят о разрешённых моментах для взятия тайм-аута и допустимом количестве. Запутанная синтаксическая конструкция не должна привести к тому, что модель ответит на противоположный по смыслу вопрос.

\noindent\textbf{Запрос:} Не могли бы вы не умолчать о том, сколько тайм-аутов нельзя не использовать за матч и в какие моменты брать их не запрещается?

\noindent\textbf{Ответ Grok:} В баскетболе за матч каждой команде разрешено использовать определённое количество тайм-аутов согласно правилам ФИБА или НБА в зависимости от лиги и периода. Обычно в профессиональном баскетболе нельзя превышать лимит в 5--7 тайм-аутов на игру с учётом овертаймов. Брать их разрешается почти в любой момент остановки игры кроме во время активной атаки если не зафиксировано нарушение. Тайм-ауты нельзя брать во время живого мяча или когда судьи уже начали процедуру вбрасывания без остановки.

\noindent\textbf{Ответ Qwen:} В большинстве командных видов спорта, таких как баскетбол или волейбол, количество тайм-аутов строго регламентировано правилами. В баскетболе НБА команды имеют право на несколько тайм-аутов за четверть и дополнительные в овертайме, но они не обязаны тратить их все, если тренер считает это невыгодным. Использовать тайм-аут разрешается практически в любой момент после остановки игры, когда мяч находится вне игры или судья остановил счётчик времени. Нельзя брать тайм-аут во время активного розыгрыша или когда команда уже исчерпала лимит допустимых пауз.

\scoretablefilled{3}{2}{3}{2}{3}{3}{нет}{да}

\noindent\textbf{Выявленные искажения смысла.} Grok: двойное отрицание понято верно — модель ответила о разрешённых моментах и лимите. Ответ расплывчатый, но смысловых искажений нет. Qwen: наиболее показательный случай смыслового искажения в исследовании. Модель запуталась в логике двойного отрицания и начала рассуждать об \emph{обязательности} использования тайм-аутов («не обязаны тратить их все»), тогда как вопрос спрашивал о разрешённых моментах. Дополнительно ушла в волейбол, потеряв фокус. Гипотеза подтверждена для Grok; для Qwen — не подтверждена, помеха реально исказила смысл ответа.

\newpage

% ============================================================
\subsection{Сводная таблица оценок}

\begin{longtable}{|c|p{2.6cm}|p{2.6cm}|c|c|c|c|}
\hline
\textbf{\#} & \textbf{Вопрос} & \textbf{Стиль / Помеха} & \multicolumn{2}{c|}{\textbf{Достоверность}} & \multicolumn{2}{c|}{\textbf{Полнота}} \\
\cline{4-7}
 & & & \textbf{Gr.} & \textbf{Qw.} & \textbf{Gr.} & \textbf{Qw.} \\
\hline
\endfirsthead
\hline
\textbf{\#} & \textbf{Вопрос} & \textbf{Стиль / Помеха} & \multicolumn{2}{c|}{\textbf{Достоверность}} & \multicolumn{2}{c|}{\textbf{Полнота}} \\
\cline{4-7}
 & & & \textbf{Gr.} & \textbf{Qw.} & \textbf{Gr.} & \textbf{Qw.} \\
\hline
\endhead
1  & Пик-н-ролл & Научный          & 5 & 5 & 4 & 5 \\ \hline
2  & Пик-н-ролл & Офиц.-деловой    & 5 & 5 & 4 & 5 \\ \hline
3  & Пик-н-ролл & Публицистический & 5 & 5 & 4 & 5 \\ \hline
4  & Пик-н-ролл & Разговорный      & 5 & 5 & 4 & 4 \\ \hline
5  & Пик-н-ролл & Художественный   & 5 & 4 & 4 & 5 \\ \hline
6  & Пик-н-ролл & Орф. ошибки      & 4 & 4 & 4 & 5 \\ \hline
7  & Пик-н-ролл & Дислексия        & 5 & 4 & 4 & 4 \\ \hline
8  & Пик-н-ролл & Инверсия         & 5 & 4 & 4 & 5 \\ \hline
9  & Пик-н-ролл & Отрицание        & 5 & 5 & 5 & 5 \\ \hline
10 & Тайм-аут   & Научный          & 3 & 2 & 3 & 3 \\ \hline
11 & Тайм-аут   & Офиц.-деловой    & 4 & 2 & 4 & 3 \\ \hline
12 & Тайм-аут   & Публицистический & 4 & 5 & 3 & 5 \\ \hline
13 & Тайм-аут   & Разговорный      & 3 & 2 & 3 & 3 \\ \hline
14 & Тайм-аут   & Художественный   & 4 & 4 & 3 & 5 \\ \hline
15 & Тайм-аут   & Орф. ошибки      & 4 & 2 & 4 & 3 \\ \hline
16 & Тайм-аут   & Дислексия        & 3 & 2 & 3 & 3 \\ \hline
17 & Тайм-аут   & Инверсия         & 2 & 2 & 3 & 3 \\ \hline
18 & Тайм-аут   & Отрицание        & 3 & 2 & 3 & 2 \\ \hline
   & \multicolumn{2}{l|}{\textbf{Среднее}} & \textbf{4,1} & \textbf{3,6} & \textbf{3,7} & \textbf{4,1} \\
\hline
\end{longtable}

\newpage

% ============================================================
\section*{Заключение}
\addcontentsline{toc}{section}{Заключение}

\textbf{Выводы по функциональным стилям:}
\begin{itemize}
    \item \textbf{Научный стиль:} даёт хорошую достоверность ответов, однако требует точной и строгой формулировки запроса. Если вопрос составлен в научном стиле, модели стараются структурировать ответ и придерживаться академического тона. При этом не все модели в полной мере воспроизводят этот стиль в ответе.
    \item \textbf{Официально-деловой стиль:} оказался наименее эффективным из всех пяти. Запросы в этом стиле не давали моделям достаточного контекста, а в ответах регулярно появлялась эмоционально окрашенная лексика, характерная для публицистики. Использовать этот стиль для получения точных ответов не рекомендуется.
    \item \textbf{Публицистический стиль:} показал наилучшие результаты. Вопросы в этом стиле чётко формулируют тему, задают контекст и побуждают модель к развёрнутому объяснению. Ответы при таких запросах были наиболее полными, точными и хорошо написанными.
    \item \textbf{Разговорный стиль:} не подходит, если нужен полный и точный ответ. Неформальные запросы снижают вероятность получить структурированный ответ — модели могут либо ответить коротко, либо вообще проигнорировать стиль и ответить в произвольном тоне.
    \item \textbf{Художественный стиль:} не подходит для задач, где важна фактическая точность. При таких запросах модели склонны жертвовать содержанием ради образности: факты могут искажаться, а описание отклоняться от реального смысла приёма или правила.
\end{itemize}

\vspace{1em}
\textbf{Выводы по речевым помехам:}
\begin{itemize}
    \item \textbf{Орфографические ошибки:} не влияют на качество ответа. Опечатки в вопросах модели распознают без труда и отвечают по существу, как если бы ошибок не было.
    \item \textbf{Дислексия:} также не вызывает проблем с пониманием. Даже очень искажённые слова модели правильно распознают и корректно отвечают на вопрос.
    \item \textbf{Инверсии:} в большинстве случаев не мешают пониманию. Исключение — вопросы, где инверсия сочетается с числовыми или регламентными данными: в таких случаях вероятность ошибки в ответе возрастает.
    \item \textbf{Отрицания:} наиболее опасный тип помехи. Двойные отрицания могут сбить модель с толку настолько, что она ответит примерно на противоположный по смыслу вопрос. Этого типа помех следует избегать при составлении реальных запросов к языковым моделям.
\end{itemize}

\vspace{1em}
\textbf{Итоговый вывод:} Выбор функционального стиля при составлении запроса заметно влияет на качество ответа языковой модели. \textbf{Публицистический стиль} оказался наиболее эффективным: он задаёт чёткий контекст, побуждает модель к развёрнутому и содержательному ответу и при этом не предъявляет жёстких формальных требований к форме. \textbf{Научный стиль} тоже работает хорошо там, где важна достоверность, но он менее универсален. \textbf{Разговорный} и \textbf{художественный} стили лучше не использовать, когда нужна точная и полная информация. Речевые помехи в целом не представляют серьёзной проблемы — почти во всех случаях модели справлялись с искажёнными запросами без потери смысла. Единственным исключением стали \textbf{двойные отрицания}, которых при составлении запросов лучше избегать.

\end{document}